\documentclass[11pt]{article}

\usepackage[margin=1.1in]{geometry}
\usepackage{amsmath,amssymb,amsthm,mathtools}
\usepackage{booktabs}
\usepackage{float}
\usepackage{array}
\usepackage{enumitem}
\usepackage{hyperref}

\hypersetup{
  colorlinks=true,
  linkcolor=blue,
  citecolor=blue,
  urlcolor=blue
}

\newtheorem{theorem}{Theorem}[section]
\newtheorem{lemma}[theorem]{Lemma}
\newtheorem{corollary}[theorem]{Corollary}
\newtheorem{proposition}[theorem]{Proposition}
\theoremstyle{definition}
\newtheorem{definition}[theorem]{Definition}
\newtheorem{problem}[theorem]{Problem}
\newtheorem{remark}[theorem]{Remark}

\DeclareMathOperator{\rad}{rad}
\DeclareMathOperator{\lcm}{lcm}
\DeclareMathOperator{\ord}{ord}
\newcommand{\docversion}{20}
\newcommand{\Pplus}{P^{+}}
\newcommand{\thetafun}{\vartheta}
\newcommand{\Prad}{P_r}
\newcommand{\BrD}{B_r(D)}
\newcommand{\etaD}{\eta_D}
\newcommand{\RD}{R_D}

\title{Per-Branch Permanence and Support-Purity Failure
for Prime-Complete Consecutive Integer Products (Version \docversion)}
\author{Ken Clements}
\date{\today}

\begin{document}
\maketitle

\begin{abstract}
Let
\[
  P_r=p_r^\#=\prod_{i\le r}p_i,
\]
and call a consecutive product prime-complete of order $r$ if
\[
  \rad(m(m+1))=P_r.
\]
Every such product gives a Pell branch
\[
  x^2-Dy^2=1,
  \qquad D\mid P_r,
\]
where $D$ is the squarefree part of $m(m+1)$.  On a branch $D$, the primes in
$P_r/D$ form the complementary support: a prime-complete candidate requires
$P_r/D\mid y_n$ and also requires $y_n$ to contain no prime exceeding $p_r$.
This paper proves a per-branch permanence theorem: for every fixed squarefree
$D>1$, there is a threshold $r_0(D)$ such that for all $r\ge r_0(D)$ no Pell
term on branch $D$ can be both support-complete and $p_r$-smooth.  The proof is
unconditional and combines a support-mass lower bound with Stewart's fixed-sequence
large-prime theorem for Lucas sequence terms.

The theorem shows that no fixed Pell branch can host prime-complete products at
arbitrarily large orders.  Any infinite tail, if it exists, must therefore pass
through fresh discriminants whose largest prime factor tends to infinity.  We then
record the observed small-index diagnostics for $8\le r\le 14$: after the known
order-$8$ prime-complete hit, all pincer-surviving small-index rows are
support-complete but support-impure.  Thus the remaining Track 2 obstruction is
not support coverage; it is support purity.
\end{abstract}

\section{Introduction}

Let $p_r$ be the $r$-th prime and let
\[
  P_r=p_r^\#=\prod_{i=1}^{r}p_i.
\]
We call $m(m+1)$ \emph{prime-complete of order $r$} if
\[
  \rad(m(m+1))=P_r.
\]
Equivalently, every prime $p\le p_r$ divides $m(m+1)$ and no prime $>p_r$ divides
$m(m+1)$.

The identity
\[
  (2m+1)^2-1=4m(m+1)
\]
places prime-complete products on Pell branches.  If $D$ is the squarefree part
of $m(m+1)$, then $D\mid P_r$ and
\[
  x=2m+1,\qquad x^2-Dy^2=1
\]
for a suitable integer $y$.  Thus, for each squarefree $D>1$, one has a Pell
branch
\[
  x_n+y_n\sqrt D=(x_1+y_1\sqrt D)^n,
  \qquad n\ge 1,
\]
where $x_1+y_1\sqrt D$ is the fundamental solution of $x^2-Dy^2=1$.

The older form of the permanence argument can be summarized as follows:
\[
  \text{fixed }D \quad \Longrightarrow \quad \text{eventual branch death.}
\]
That statement remains correct and is proved below.  The new diagnostic point is
sharper.  On branch $D$, define the complementary product
\[
  B_r(D)=\frac{P_r}{D}.
\]
A prime-complete candidate must satisfy
\[
  B_r(D)\mid y_n
\]
and also
\[
  \Pplus(y_n)\le p_r.
\]
The first condition is \emph{support completion}; the second is \emph{support
purity}.  The computations reported here show that after the known order-$8$
solution, the small-index residue is not failing because required primes are
missing.  It is failing because support-complete terms acquire outside primes.

Thus the conceptual picture becomes
\[
  \boxed{\text{after }r=8,\ \text{support survives, but purity dies.}}
\]
The per-branch theorem explains why old branches cannot persist forever.  The
support-purity diagnostics explain what happens microscopically in the first
orders after the last known hit.

\section{Notation and branch conditions}

For $N\ge 2$, write $\Pplus(N)$ for the largest prime divisor of $N$, and set
$\Pplus(1)=1$ by convention.  Let
\[
  \thetafun(x)=\sum_{p\le x}\log p.
\]
Then
\[
  \log P_r=\thetafun(p_r).
\]

\begin{definition}[Prime-complete product]
A product $m(m+1)$ is \emph{prime-complete of order $r$} if
\[
  \rad(m(m+1))=P_r.
\]
\end{definition}

\begin{definition}[Pell branch]
Let $D>1$ be squarefree.  Let
\[
  \eta_D=x_1+y_1\sqrt D>1
\]
be the fundamental solution of
\[
  x^2-Dy^2=1.
\]
The positive solutions are indexed by
\[
  x_n+y_n\sqrt D=\eta_D^n,
  \qquad n\ge1.
\]
Set
\[
  R_D=\log \eta_D.
\]
\end{definition}

\begin{definition}[Support completion and support purity]
Suppose $D\mid P_r$.  Define
\[
  B_r(D)=\frac{P_r}{D}.
\]
A Pell term $y_n$ on branch $D$ is \emph{$r$-support-complete} if
\[
  B_r(D)\mid y_n.
\]
It is \emph{$p_r$-smooth}, or \emph{support-pure at order $r$}, if
\[
  \Pplus(y_n)\le p_r.
\]
It is a \emph{support-purity failure} if
\[
  B_r(D)\mid y_n
  \qquad\text{but}\qquad
  \Pplus(y_n)>p_r.
\]
\end{definition}

A genuine prime-complete product must produce both support completion and support
purity on its Pell branch.  This is the basic branch attachment.

\section{Prime-complete products give support-complete, support-pure Pell terms}

\begin{lemma}[Pell attachment]\label{lem:pell-attach}
Let $m(m+1)$ be prime-complete of order $r$, and let $D$ be the squarefree part
of $m(m+1)$.  Then $D\mid P_r$, and there exist integers $x,y$ with
\[
  x=2m+1,
  \qquad
  x^2-Dy^2=1.
\]
Moreover, $(x,y)=(x_n,y_n)$ for some $n\ge1$ on branch $D$.
\end{lemma}

\begin{proof}
Since $D$ is the squarefree part of $m(m+1)$, there is an integer $z\ge1$ such
that
\[
  m(m+1)=Dz^2.
\]
Prime-completeness implies every prime divisor of $D$ lies among
$p_1,\ldots,p_r$, so $D\mid P_r$.  Now
\[
  (2m+1)^2-1=4m(m+1)=4Dz^2.
\]
With $x=2m+1$ and $y=2z$, this gives $x^2-Dy^2=1$.  All positive solutions of
Pell's equation are powers of the fundamental unit, so $(x,y)=(x_n,y_n)$ for
some $n\ge1$.
\end{proof}

\begin{lemma}[Complementary support and purity]\label{lem:support}
Let $m(m+1)$ be prime-complete of order $r$, and let $D$ be the squarefree part
of $m(m+1)$.  Let $(x_n,y_n)$ be the associated Pell solution.  Then
\[
  B_r(D)=\frac{P_r}{D}\mid y_n
\]
and
\[
  \Pplus(y_n)\le p_r.
\]
\end{lemma}

\begin{proof}
Write
\[
  m(m+1)=Dz^2,
  \qquad
  y_n=2z.
\]
Let $p\le p_r$ with $p\nmid D$.  Since $m(m+1)$ is prime-complete, $p$ divides
$m(m+1)$.  Since $p\nmid D$, the exponent of $p$ in $Dz^2$ is even and positive,
hence at least $2$.  Therefore $p\mid z$, and hence $p\mid y_n$.  This holds for
every prime divisor of $P_r/D$, so $P_r/D\mid y_n$.

Finally, $y_n=2z$ has no prime divisor outside the prime divisors of $m(m+1)$,
so $\Pplus(y_n)\le p_r$.
\end{proof}

\section{Reduction-theorem context}

The per-branch permanence theorem is independent of the small-index reduction,
but the two results frame the same residual problem.  The reduction theorem says
that any tail prime-complete product must occur in a small-index, support-complete
core.

\begin{theorem}[Reduction theorem, quoted form]\label{thm:reduction-context}
Let $r$ be in the tail range, and suppose a prime-complete product of order $r$
exists.  Then for some squarefree $D\mid P_r$ and some Pell index $n$ on branch
$D$, the associated term $y_n$ is support-complete and $p_r$-smooth, and the
index satisfies
\[
  n\le p_r+1.
\]
Moreover the complementary rank-lcm condition holds:
\[
  R_r(D)\mid n,
\]
where $R_r(D)$ is the least common multiple of the ranks of apparition of the
complementary primes in the branch sequence.
\end{theorem}

\begin{remark}
This theorem is not reproved here.  It uses primitive divisors to cap the Pell
index and rank-of-apparition divisibility to force the complementary primes into
a common small index.  The present paper proves a perpendicular statement: for a
fixed branch, support completion eventually forces indices so large that
Stewart's fixed-sequence theorem makes $p_r$-smoothness impossible.
\end{remark}

\section{The fixed-branch large-prime input}

We use the following fixed-sequence form of Stewart's theorem on greatest prime
factors of Lucas and Lehmer sequence terms.

\begin{theorem}[Stewart, fixed-sequence form]\label{thm:stewart}
Let $(U_n)_{n\ge1}$ be a nondegenerate Lucas sequence with real roots.  Then
there exists an effectively computable integer $n_0(U)$ such that, for every
$n>n_0(U)$,
\[
  \Pplus(U_n)
  >
  n\exp\!\left(\frac{\log n}{104\log\log n}\right).
\]
\end{theorem}

\begin{remark}
The constant $104$ is not important for the permanence argument.  What matters is
that
\[
  \exp\!\left(\frac{\log n}{104\log\log n}\right)\to\infty.
\]
The theorem is used one fixed branch at a time, so the threshold may depend on
$D$.  No branch-uniform Stewart threshold is assumed.
\end{remark}

For branch $D$, the normalized sequence
\[
  U_n(D)=\frac{y_n}{y_1}
\]
is an integer Lucas sequence with real roots.  Since $U_n(D)\mid y_n$,
Stewart's theorem supplies a lower bound for $\Pplus(y_n)$ whenever $n$ exceeds
the branch-dependent threshold.

\section{Support mass}

The elementary mass estimate is the first jaw.  If $B_r(D)$ divides $y_n$, then
$y_n$ must be at least $B_r(D)$, which forces the index $n$ to grow with $r$.

\begin{lemma}[Support-mass lower bound]\label{lem:mass}
Fix squarefree $D>1$ and let $R_D=\log\eta_D$.  Suppose $D\mid P_r$ and
\[
  B_r(D)=P_r/D\mid y_n.
\]
Then
\[
  n\ge
  M_D(r):=
  \frac{\thetafun(p_r)+\log 2-\frac12\log D}{R_D}.
\]
In particular, $M_D(r)\to\infty$ as $r\to\infty$.
\end{lemma}

\begin{proof}
The Pell $y$-coordinate satisfies
\[
  y_n=\frac{\eta_D^n-\eta_D^{-n}}{2\sqrt D}
  <
  \frac{\eta_D^n}{2\sqrt D}.
\]
If $B_r(D)\mid y_n$, then $y_n\ge B_r(D)=P_r/D$.  Hence
\[
  \frac{P_r}{D}
  \le y_n
  <
  \frac{\eta_D^n}{2\sqrt D}.
\]
Therefore
\[
  \eta_D^n>\frac{2P_r}{\sqrt D}.
\]
Taking logarithms gives
\[
  nR_D>\thetafun(p_r)+\log 2-\frac12\log D.
\]
Replacing $>$ by $\ge$ in the displayed lower bound is harmless; one may insert a
ceiling if an integral bound is desired.  Since $D$ and $R_D$ are fixed while
$\thetafun(p_r)\to\infty$, the final assertion follows.
\end{proof}

\section{Per-branch permanence}

\begin{theorem}[Per-Branch Permanence]\label{thm:permanence}
Fix a squarefree integer $D>1$.  There exists an integer $r_0(D)$ such that for
every $r\ge r_0(D)$ there is no index $n\ge1$ for which both
\[
  B_r(D)=P_r/D\mid y_n
\]
and
\[
  \Pplus(y_n)\le p_r
\]
hold.  Equivalently, for all sufficiently large $r$, branch $D$ carries no term
that is simultaneously support-complete and support-pure at order $r$.
\end{theorem}

\begin{proof}
Fix $D$.  Let
\[
  \eta_D=x_1+y_1\sqrt D,
  \qquad
  R_D=\log\eta_D,
\]
and let
\[
  U_n(D)=\frac{y_n}{y_1}
\]
be the associated Lucas sequence.  By Theorem~\ref{thm:stewart}, there exists a
branch-dependent threshold $n_0(D)$ such that, for all $n>n_0(D)$,
\[
  \Pplus(U_n(D))
  >
  n\exp\!\left(\frac{\log n}{104\log\log n}\right).
\]
Since $U_n(D)\mid y_n$, the same lower bound applies to $\Pplus(y_n)$.

Let
\[
  M_D(r)=\frac{\thetafun(p_r)+\log 2-\frac12\log D}{R_D}.
\]
By Lemma~\ref{lem:mass}, any index $n$ satisfying $B_r(D)\mid y_n$ must have
\[
  n\ge M_D(r).
\]
By the prime number theorem in the form
\[
  \thetafun(p_r)\sim p_r,
\]
we have
\[
  M_D(r)\sim \frac{p_r}{R_D}.
\]
In particular $M_D(r)\to\infty$.  Since
\[
  \exp\!\left(\frac{\log M_D(r)}{104\log\log M_D(r)}\right)\to\infty,
\]
and $R_D$ is fixed, we may choose $r_0(D)$ sufficiently large that for all
$r\ge r_0(D)$:
\begin{enumerate}[label=(\roman*)]
  \item $D\mid P_r$;
  \item $M_D(r)>n_0(D)$;
  \item the function
  \[
    t\mapsto t\exp\!\left(\frac{\log t}{104\log\log t}\right)
  \]
  is increasing for $t\ge M_D(r)$;
  \item
  \[
    M_D(r)\exp\!\left(\frac{\log M_D(r)}{104\log\log M_D(r)}\right)>p_r.
  \]
\end{enumerate}
Now suppose $r\ge r_0(D)$ and $B_r(D)\mid y_n$.  Then
\[
  n\ge M_D(r)>n_0(D).
\]
Stewart's lower bound and the monotonicity condition give
\[
  \Pplus(y_n)
  \ge
  \Pplus(U_n(D))
  >
  n\exp\!\left(\frac{\log n}{104\log\log n}\right)
  \ge
  M_D(r)\exp\!\left(\frac{\log M_D(r)}{104\log\log M_D(r)}\right)
  >p_r.
\]
Thus $y_n$ is not $p_r$-smooth.  Hence the simultaneous requirements
$B_r(D)\mid y_n$ and $\Pplus(y_n)\le p_r$ are impossible for all $r\ge r_0(D)$.
\end{proof}

\begin{corollary}[No fixed branch supports infinitely many prime-complete products]
\label{cor:fixedbranch}
Fix squarefree $D>1$.  There are only finitely many orders $r$ for which a
prime-complete product of order $r$ has squarefree part $D$.
\end{corollary}

\begin{proof}
A prime-complete product on branch $D$ gives, by Lemma~\ref{lem:support}, an
index $n$ with
\[
  B_r(D)\mid y_n
  \qquad\text{and}\qquad
  \Pplus(y_n)\le p_r.
\]
Theorem~\ref{thm:permanence} rules this out for all sufficiently large $r$.
\end{proof}

\begin{corollary}[Fresh-branch corollary]\label{cor:fresh}
Let $S$ be a fixed finite set of squarefree integers $D>1$.  There exists
$r_0(S)$ such that, for all $r\ge r_0(S)$, no prime-complete product of order
$r$ has squarefree part in $S$.

In particular, if prime-complete products existed at arbitrarily large orders,
their squarefree parts $D_r$ would satisfy
\[
  \Pplus(D_r)\to\infty.
\]
\end{corollary}

\begin{proof}
For finite $S$, take
\[
  r_0(S)=\max_{D\in S}r_0(D).
\]
The first assertion follows from Corollary~\ref{cor:fixedbranch}.  For the
second, if $\Pplus(D_r)$ did not tend to infinity, then infinitely many $D_r$
would lie in a fixed finite set of squarefree integers composed only of primes up
to some fixed bound, contradicting the first assertion.
\end{proof}

\section{Support-purity diagnostics}

This section is diagnostic, not part of the proof of Theorem~\ref{thm:permanence}.
It records what happens in the exhaustive branch and small-index computations for
$8\le r\le 14$.  The branch pincer leaves a finite small-index residue.  For each
row in that residue the computation records whether the complementary support is
covered, whether $y_n$ is $p_r$-smooth, and whether the row produces a direct
prime-complete hit.

\begin{table}[H]
\centering
\small
\begin{tabular}{rrrrrrrrr}
\toprule
$r$ & rows & branches & supp & $y$-sm & both & pc & purfail & min left \\
\midrule
8  & 136 & 52 & 136 & 1 & 1 & 1 & 135 & 2.253 \\
9  & 145 & 57 & 145 & 0 & 0 & 0 & 145 & 2.029 \\
10 & 110 & 52 & 110 & 0 & 0 & 0 & 110 & 5.150 \\
11 & 91  & 47 & 91  & 0 & 0 & 0 & 91  & 12.867 \\
12 & 50  & 30 & 50  & 0 & 0 & 0 & 50  & 12.824 \\
13 & 41  & 26 & 41  & 0 & 0 & 0 & 41  & 22.435 \\
14 & 11  & 9  & 11  & 0 & 0 & 0 & 11  & 81.862 \\
\bottomrule
\end{tabular}
\caption{Support-purity diagnostic for $8\le r\le 14$.  Here ``supp'' means
support-complete, ``$y$-sm'' means $y_n$ is $p_r$-smooth, ``both'' means both
support-complete and $p_r$-smooth, and ``purfail'' means support-complete but not
$p_r$-smooth.  The final column is $\log_{10}$ of the smallest leftover after
stripping primes $\le p_r$ from $y_n$ among support-purity failures.}
\label{tab:purity-order}
\end{table}

Table~\ref{tab:purity-order} shows the central point.  At $r=8$ there is exactly
one row that is support-complete, $p_r$-smooth, and prime-complete.  For every
$r=9,\ldots,14$, all small-index rows attached to pincer survivors are
support-complete, but none has $p_r$-smooth $y_n$.

Thus after $r=8$ the obstruction is not missing support.  It is the presence of
outside primes:
\[
  B_r(D)\mid y_n
  \qquad\text{but}\qquad
  \Pplus(y_n)>p_r.
\]
Moreover, the closest support-purity failures move rapidly away from purity.  By
$r=14$, even the closest candidate has an outside leftover of size about
$10^{81.862}$ after all primes $\le p_{14}$ have been stripped from $y_n$.

The closest global failures in the computation begin as follows.
\begin{table}[H]
\centering
\small
\begin{tabular}{rrrrrlr}
\toprule
rank & $r$ & mask & $D$ & $n$ & origin & $\log_{10}$ leftover \\
\midrule
1 & 9  & 0x16b & 213486 & 2 & newbit/dead-parent & 2.029 \\
2 & 8  & 0x67  & 6630   & 2 & base & 2.253 \\
3 & 9  & 0x154 & 21505  & 3 & newbit/live-parent & 3.303 \\
4 & 9  & 0xa3  & 1482   & 6 & oldmask/live-parent & 4.505 \\
5 & 9  & 0xb5  & 27170  & 4 & oldmask/live-parent & 4.619 \\
6 & 8  & 0x53  & 1122   & 6 & base & 4.966 \\
7 & 10 & 0x2d0 & 103037 & 2 & newbit/dead-parent & 5.150 \\
\bottomrule
\end{tabular}
\caption{Closest support-purity failures among the diagnostic rows.  The nearest
post-$8$ miss occurs at $r=9$, $D=213486$, $n=2$, and is already impure.}
\label{tab:closest}
\end{table}

The origin labels separate inherited old-mask survivors from genuinely new-bit
discriminant branches.  The closest post-$8$ failures are often fresh-branch
leakage: a new discriminant branch survives the first pincer but fails purity
immediately.

Two quantitative features of this residue are worth recording, both diagnostic.
First, the closest-miss margin grows geometrically.  Writing $\lambda(r)$ for the
$\log_{10}$ leftover of the nearest support-purity failure at order $r$
(Table~\ref{tab:purity-order}, final column), the values
$2.03,\,5.15,\,12.87,\,12.82,\,22.43,\,81.86$ for $r=9,\ldots,14$ are fit by
$\lambda(r)\approx c\cdot\rho^{r}$ with $\rho\approx 1.9$: even the single most
threatening candidate at each order recedes from purity at a geometric rate.  We
attach the standing caveat that $r=14$ rests on only $11$ rows from $9$ branches,
so the rate constant is known only to within perhaps $\pm20\%$; the qualitative
divergence, however, is uniform across the range.

Second, this purity divergence is the microscopic complement of the branch-level
survival collapse.  The branch-lineage census over the same range shows that the
fraction of newly created discriminant branches that survive the first pincer
decays geometrically, at rate $\approx 0.37^{r}$; the present diagnostic shows
that those few survivors then fail purity by a margin growing at $\approx 1.9^{r}$.
The small-index residue is thus emptied from two sides at once: fewer branches
survive support, and the survivors fall ever further from purity.  Both are
average-case statements about the residue as a whole.

Neither statement touches the fundamental case $n=1$.  The single support-pure
row in the entire scan is the order-$8$ solution, which occurs at $n=1$; and $n=1$
is precisely where the size mechanism behind the purity divergence has no purchase
(there is no index to grow and no higher Lucas term to force an outside prime).
The diagnostics therefore corroborate purity failure for $n\ge 2$ with a growing
margin, and leave the $n=1$ fundamental-solution case as the irreducible residue,
consistent with the residual problem stated in Section~\ref{sec:residual}.

\section{Relation to the wall $A_r$ versus $L_r$}

Let $L_r$ denote the largest $m$ such that $m$ and $m+1$ are both $p_r$-smooth,
and let $A_r$ denote the least admissible CRT floor value at which
\[
  P_r\mid m(m+1)
\]
can occur at product scale.  A prime-complete product of order $r$ must lie in
the coarse interval
\[
  A_r\le m\le L_r.
\]
Therefore
\[
  A_r>L_r
  \qquad\Longrightarrow\qquad
  \text{no prime-complete product of order }r.
\]

The wall first closes in the available data at $r=14$: at $r=13$ one still has
$A_{13}<L_{13}$, while at $r=14$ one has $A_{14}>L_{14}$.  The support-purity
diagnostics explain this crossover microscopically.

At $r=13$, the coarse wall is still open, so the interval argument alone does not
exclude a prime-complete product.  But the branch-level residue already has
\[
  41\text{ rows},\qquad 41\text{ support-complete},\qquad 0\text{ }p_{13}\text{-smooth}.
\]
The closest outside leftover has $\log_{10}$ about $22.435$.  Thus the candidates
are dead by support-purity failure before the global wall is visible.

At $r=14$, both views agree.  The wall has closed, and the branch-level residue is
more impure still:
\[
  11\text{ rows},\qquad 11\text{ support-complete},\qquad 0\text{ }p_{14}\text{-smooth},
\]
with closest outside leftover about $10^{81.862}$.

Thus the wall is the first global interval closure, while support-purity failure
is the local branch mechanism already visible before the wall closes.

\section{The rewritten residual problem}\label{sec:residual}

Theorem~\ref{thm:permanence} proves that old fixed branches die.  The reduction
theorem confines any tail counterexample to small-index support-complete terms.
The diagnostics indicate that the remaining obstruction is support purity.  This
motivates the following residual problem.

\begin{problem}[Small-Index Support-Purity Failure Problem]\label{prob:purity}
Prove that, for all sufficiently large $r$, there is no squarefree $D\mid P_r$
and no index
\[
  1\le n\le p_r+1
\]
such that
\[
  B_r(D)=P_r/D\mid y_n(D)
\]
and
\[
  \Pplus(y_n(D))\le p_r.
\]
Equivalently, prove that in the small-index range forced by the reduction
theorem,
\[
  P_r/D\mid y_n(D)
  \qquad\Longrightarrow\qquad
  \Pplus(y_n(D))>p_r
\]
for all sufficiently large $r$.
\end{problem}

Problem~\ref{prob:purity} is strictly sharper than a generic smooth-pair ceiling.
It does not ask whether all consecutive $p_r$-smooth pairs are bounded.  It asks
whether the specific support-complete Pell terms forced by the branch reduction
can remain support-pure.  The data suggests that support completion and support
purity coincide for the last time at order $8$.

\section{Interpretation and limitations}

The Per-Branch Permanence Theorem is a structural ratchet.  It proves that no
fixed $D$ can support prime-complete products forever.  Thus a hypothetical
infinite tail cannot live on a stable finite set of discriminants; it must pass
through fresh branches containing larger and larger primes.

This is not yet a global unconditional termination proof.  The threshold
$r_0(D)$ depends on $D$, and the explicit constants in Stewart's theorem are far
too large to produce a useful uniform bound across all $D\mid P_r$.  The young
branches remain the hard part.

The new support-purity diagnostics refine the residual problem.  The issue is not
that support-complete small-index rows disappear immediately.  In the observed
range, they are abundant among pincer survivors.  The issue is that after $r=8$
they are never support-pure.  The branch-level slogan is therefore:
\[
  \boxed{\text{support completion persists; support purity fails.}}
\]

\begin{thebibliography}{9}

\bibitem{Stewart2013}
C. L. Stewart,
\emph{On divisors of Lucas and Lehmer numbers},
Acta Mathematica \textbf{211} (2013), 291--314.

\bibitem{BHV2001}
Y. Bilu, G. Hanrot, and P. M. Voutier,
\emph{Existence of primitive divisors of Lucas and Lehmer numbers},
Journal f\"ur die reine und angewandte Mathematik \textbf{539} (2001), 75--122.

\bibitem{Lehmer1964}
D. H. Lehmer,
\emph{On a problem of St\o rmer},
Illinois Journal of Mathematics \textbf{8} (1964), 57--79.

\bibitem{LucaNajman2010}
F. Luca and F. Najman,
\emph{On the largest prime factor of $x^2-1$},
arXiv:1005.1533.

\bibitem{Najman2010}
F. Najman,
\emph{Smooth values of quadratic polynomials},
arXiv:1005.1535.

\end{thebibliography}

\end{document}
